\section{Scalar soft truncation}

For $t>0$, an integer $m\geq0$, and $x\geq0$, the scalar soft truncation of
Arad--Landau--Vazirani--Vidick is the following finite logarithmic series
~\cite[lines 1152--1173]{Arad2016RigorousRG}.

\begin{definition}[Scalar soft truncation]
    \label{def:scalar_soft_truncation}
    \lean{Real.softTruncation}
    \leanok
    \begin{align}
        h_{m,t}(x)
            &=t\sum_{j=1}^{m}\frac{{(1-e^{-x/t})}^{j}}{j}.
        \label{eq:scalar_soft_truncation_definition}
    \end{align}
\end{definition}

\begin{theorem}[Scalar soft-truncation bounds]
    \label{thm:scalar_soft_truncation_bounds}
    \lean{Real.softTruncation_bounds}
    \leanok
    If $t>0$, $m\geq1$, and $x\geq0$, then
    \begin{align}
        0\leq h_{m,t}(x)&\leq x,
        \label{eq:scalar_soft_truncation_order}\\
        x-h_{m,t}(x)&\leq\frac{t}{m}{\left(\frac{x}{t}\right)}^{m},
        \label{eq:scalar_soft_truncation_remainder}\\
        h_{m,t}(x)&\leq t(1+\log m).
        \label{eq:scalar_soft_truncation_log_bound}
    \end{align}
    The final estimate corrects the printed $t\log m$ bound, which fails at
    $m=1$~\cite{gap:alvv16_soft_truncation_log_bound}.
\end{theorem}
\begin{proof}\leanok
    Put $y=x/t$, $a(y)=1-e^{-y}$, and
    $g_m(y)=\sum_{j=1}^{m}{a(y)}^{j}/j$. The finite geometric sum gives
    $g_m'(y)=1-{a(y)}^{m}$. Since $0\leq a(y)\leq\min\{1,y\}$ for $y\geq0$,
    the functions $y-g_m(y)$ and $g_m(y)-y+y^m/m$ are nondecreasing and
    vanish at zero. This proves~(\ref{eq:scalar_soft_truncation_order}) and
    (\ref{eq:scalar_soft_truncation_remainder}). Finally,
    ${a(y)}^{j}\leq1$, so the harmonic estimate
    $\sum_{j=1}^{m}1/j\leq1+\log m$ proves
    (\ref{eq:scalar_soft_truncation_log_bound}).
\end{proof}


\begin{lemma}[Geometric remainder]
    \label{lem:scalar_soft_truncation_geometric}
    \lean{Real.softTruncation_remainder_le_geometric}
    \leanok
    If $t>0$, $m\geq0$ is an integer, and $x\geq0$, then
    \begin{align}
        x-h_{m,t}(x)\leq x(1-e^{-x/t})^m.
        \label{eq:scalar_soft_truncation_geometric}
    \end{align}
\end{lemma}
\begin{proof}\leanok
    The derivative of the remainder at $u\geq0$ is $(1-e^{-u/t})^m$.
    On $[0,x]$ it is bounded above by $(1-e^{-x/t})^m$, and the
    remainder vanishes at zero. The mean value inequality gives the result.
\end{proof}

\begin{theorem}[Exponential weight for the remainder]
    \label{thm:scalar_soft_truncation_exp_weight}
    \lean{Real.softTruncation_remainder_le_exp}
    \leanok
    If $m\geq1$ is an integer, $0<\theta\leq1/2$, and $x\geq0$, then
    \begin{align}
        x-h_{m,m}(x)
        \leq\frac2\theta\exp\!\left(\theta x-\frac{\theta m}{2}\right).
        \label{eq:scalar_soft_truncation_exp_weight}
    \end{align}
\end{theorem}
\begin{proof}\leanok
    The inequality $z\leq e^{sz}/s$ holds for $s>0$. If $x\leq m$,
    then
    \begin{align}
        1-e^{-x/m}\leq1-e^{-1}\leq\frac34\leq e^{-\theta/2}.
    \end{align}
    Applying~(\ref{eq:scalar_soft_truncation_geometric}) gives
    $x-h_{m,m}(x)\leq\theta^{-1}e^{\theta x-\theta m/2}$.
    If $x\geq m$, the nonnegativity of the truncation instead gives
    \begin{align}
        x-h_{m,m}(x)\leq x
        \leq\frac2\theta e^{\theta x/2}
        \leq\frac2\theta e^{\theta x-\theta m/2}.
    \end{align}
\end{proof}

For a probability spectral measure with exponential moment at most $M$,
integration gives a mean remainder at most
$2M\theta^{-1}e^{-\theta m/2}$. Application to a regional Hamiltonian
requires a bound on that physical exponential moment.
